% Petits développements autour des estimateurs statistiques
% S. Jaubert. Première version : 15 décembre 2019. Version corrigée : 30 octobre 2026.
% Compilation : xelatex (deux passes)
\documentclass[11pt,a4paper]{article}
\usepackage{polyglossia}
\setmainlanguage{french}
\usepackage{fontspec}
\usepackage{amsmath,amssymb,amsthm}
\usepackage[a4paper,margin=3.2cm,headheight=14pt]{geometry}
\usepackage{xcolor}
\usepackage{fancyhdr}
\usepackage[most]{tcolorbox}
\usepackage{enumitem}
\usepackage{booktabs}
\usepackage[colorlinks=true,linkcolor=black,urlcolor=blue!60!black]{hyperref}

\pagestyle{fancy}\fancyhf{}
\renewcommand{\headrulewidth}{0.4pt}
\renewcommand{\footrulewidth}{0.4pt}
\lhead{\small Enseignement scientifique · statistique}
\lfoot{\LaTeX\ \scriptsize S. Jaubert · Estimation}
\cfoot{\thepage}
\rfoot{\scriptsize Version corrigée, octobre 2026}

\newtheoremstyle{sj}{6pt}{6pt}{\itshape}{}{\bfseries}{.}{0.5em}{}
\theoremstyle{sj}
\newtheorem*{prop}{Propriété}
\newtheorem*{thm}{Théorème}
\theoremstyle{remark}
\newtheorem*{rem}{Remarque}

\newcommand{\E}{\mathbb{E}}
\newcommand{\V}{\mathbb{V}}
\renewcommand{\P}{\mathbb{P}}
\newcommand{\N}{\mathcal{N}}
\newcommand{\Xb}{\overline{X}}
\newcommand{\dd}{\,\mathrm{d}}

\newtcolorbox{code}{colback=black!85,colframe=black!85,coltext=white,arc=2pt,left=6pt,right=6pt,top=4pt,bottom=4pt,fontupper=\small\ttfamily}
\newtcolorbox{note}[1][]{colback=black!3,colframe=black!40,boxrule=0.5pt,arc=2pt,left=6pt,right=6pt,top=4pt,bottom=4pt,fonttitle=\bfseries\small,title={#1}}

\setlength{\parindent}{0pt}
\setlength{\parskip}{6pt}
\setlength{\emergencystretch}{2em}

\begin{document}

{\raggedright\Huge\bfseries Petits développements autour\\[4pt] des estimateurs statistiques\par}

\bigskip
Il est très rare de trouver une explication claire des différentes constantes utilisées pour les calculs de capabilité. Très souvent, on nous renvoie à des tables sans autre explication. Mon propos ici est de donner les grandes lignes qui permettent le calcul de ces constantes.

Bien sûr, il faut quelques prérequis en mathématiques, en particulier être familier avec la loi du $\chi^2$ et la fonction $\Gamma$, qui prolonge la factorielle : $\Gamma(n)=(n-1)!$ pour tout entier $n\ge1$, et $\Gamma$ est définie sur les complexes sauf en $0,-1,-2,\dots$ (\href{https://fr.wikipedia.org/wiki/Fonction_gamma}{voir ici}).

\textbf{Notation.} Comme dans les tables industrielles, $\N(\mu;\sigma)$ désigne la loi normale de moyenne $\mu$ et d'\emph{écart-type} $\sigma$.

Soient $(X_i)_{1\le i\le n}$ $n$ variables aléatoires indépendantes et de même loi, d'espérance $\mu$ et d'écart-type $\sigma$. Cette hypothèse est essentielle, on le verra par la suite.

On rappelle que la variable aléatoire égale à la moyenne est
\[
\Xb=\frac{X_1+X_2+\dots+X_n}{n},
\]
et que la variable aléatoire égale à la variance des $n$ valeurs est
\[
S^2=\frac1n\sum_{i=1}^n\big(X_i-\Xb\big)^2 .
\]

\section*{1. La moyenne}

\begin{prop}
$\E(\Xb)=\mu$ et $\V(\Xb)=\dfrac{\sigma^2}{n}$, quelle que soit la loi des $X_i$.
Si de plus les $X_i$ suivent la loi $\N(\mu;\sigma)$, alors $\Xb$ suit la loi $\N\big(\mu;\frac{\sigma}{\sqrt n}\big)$.
\end{prop}

En effet : $\E(\Xb)=\frac1n\sum_{i=1}^n\E(X_i)=\frac{n\mu}{n}=\mu$, et, par indépendance,
\[
\V(\Xb)=\frac{1}{n^2}\sum_{i=1}^n\V(X_i)=\frac{n\sigma^2}{n^2}=\frac{\sigma^2}{n}.
\]
Ces deux calculs ne donnent que l'espérance et la variance de $\Xb$. Que $\Xb$ soit normale vient d'une autre propriété : une somme de variables normales \emph{indépendantes} est encore normale (stabilité de la loi normale, qui se démontre par un produit de convolution).

\begin{thm}[Théorème central limite]
Si les $X_i$ sont indépendantes et de même loi, \emph{quelconque}, d'espérance $\mu$ et d'écart-type $\sigma$ finis, alors pour $n$ « grand » la loi de $\Xb$ se rapproche de la loi $\N\big(\mu;\frac{\sigma}{\sqrt n}\big)$.
\end{thm}

C'est ce théorème qui rend les cartes de moyennes peu sensibles à la forme de la loi des pièces : la moyenne de 5 pièces est déjà bien plus « normale » qu'une pièce isolée.

\section*{2. La variance}

\subsection*{2.1 Une approche valable pour toute loi}

On écrit $S^2=\dfrac{1}{n}\sum_{i=1}^nX_i^2-\Xb^{\,2}$ (formule de König-Huygens). Alors
\begin{align*}
\E(S^2)&=\frac1n\sum_{i=1}^n\E(X_i^2)-\E\big(\Xb^{\,2}\big)=\E(X^2)-\E\big(\Xb^{\,2}\big)\\
&=\V(X)+\E(X)^2-\Big(\V(\Xb)+\E(\Xb)^2\Big)=\V(X)-\V(\Xb)\qquad\text{car }\E(\Xb)=\E(X)\\
&=\V(X)-\frac1n\V(X)=\frac{n-1}{n}\,\sigma^2 .
\end{align*}
Pour éviter le biais $\frac{n-1}{n}$, on prend pour estimateur ponctuel de la variance
\[
S_{n-1}^2=\frac{nS^2}{n-1}=\frac{\sum_{i=1}^n\big(X_i-\Xb\big)^2}{n-1},\qquad \E\big(S_{n-1}^2\big)=\sigma^2 .
\]
Ce calcul n'utilise que l'indépendance et l'existence de $\sigma$ : il vaut pour \emph{toute} loi.

\subsection*{2.2 Le cas normal : la loi du \texorpdfstring{$\chi^2$}{khi-deux}}

On suppose maintenant les $X_i$ de loi $\N(\mu;\sigma)$. On a
\[
\frac{nS^2}{\sigma^2}=\sum_{i=1}^n\left(\frac{X_i-\Xb}{\sigma}\right)^2 .
\]
Attention : les termes de cette somme ne sont ni indépendants (ils sont liés par $\sum(X_i-\Xb)=0$) ni de variance 1 (chacun a pour variance $\frac{n-1}{n}$). On ne peut donc pas conclure directement à une loi du $\chi^2(n)$. Le résultat exact est le suivant.

\begin{thm}[Fisher]
Si les $X_i$ sont indépendantes de loi $\N(\mu;\sigma)$, alors $\dfrac{nS^2}{\sigma^2}$ suit la loi du $\chi^2(n-1)$ (\href{https://fr.wikipedia.org/w/index.php?title=Loi_du_\%CF\%87\%C2\%B2&oldid=161071288}{loi du chi2}), et elle est indépendante de $\Xb$.
\end{thm}

\textbf{Idée de la preuve.} Soit $Q$ une matrice orthogonale dont la première ligne est $\big(\frac{1}{\sqrt n},\dots,\frac{1}{\sqrt n}\big)$. Le vecteur $Z=\frac{1}{\sigma}(X_1-\mu,\dots,X_n-\mu)$ a une densité qui ne dépend que de $\|z\|$ ; une rotation ne change donc pas sa loi, et les coordonnées $V_1,\dots,V_n$ de $QZ$ sont encore indépendantes de loi $\N(0;1)$. Or $V_1=\frac{\sqrt n}{\sigma}(\Xb-\mu)$ et, comme $Q$ conserve la norme,
\[
\frac{nS^2}{\sigma^2}=\frac{1}{\sigma^2}\Big(\sum(X_i-\mu)^2-n(\Xb-\mu)^2\Big)=\|QZ\|^2-V_1^2=\sum_{j=2}^nV_j^2 .
\]
C'est une somme de $n-1$ carrés de variables $\N(0;1)$ indépendantes, indépendante de $V_1$ : d'où le $\chi^2(n-1)$. On retrouve au passage l'idée de \emph{degrés de liberté} : la contrainte $\sum(X_i-\Xb)=0$ fait perdre une dimension.

Comme $\E\big(\chi^2(n-1)\big)=n-1$, on a $\E\left(\frac{nS^2}{\sigma^2}\right)=n-1$ : on retrouve l'estimateur sans biais $S_{n-1}^2$ du paragraphe 2.1.

\section*{3. L'écart-type : la constante \texorpdfstring{$c_4$}{c4}}

Nous avons un estimateur sans biais de la variance, mais il ne faut pas penser que sa racine carrée est un estimateur sans biais de l'écart-type. En effet, posons $K=\dfrac{nS^2}{\sigma^2}$, qui suit la loi du $\chi^2(n-1)$, de densité
\[
f(t)=\frac{e^{-t/2}\,t^{\frac{n-1}{2}-1}}{2^{\frac{n-1}{2}}\,\Gamma\big(\frac{n-1}{2}\big)}\qquad(t>0).
\]
Alors, avec le changement de variable $t=2u$ et la définition $\Gamma(a)=\int_0^{+\infty}u^{a-1}e^{-u}\dd u$ :
\[
\E\big(\sqrt K\big)=\int_0^{+\infty}\sqrt t\,f(t)\dd t=\frac{2^{\frac n2}\,\Gamma\big(\frac n2\big)}{2^{\frac{n-1}{2}}\,\Gamma\big(\frac{n-1}{2}\big)}=\frac{\sqrt2\,\Gamma\big(\frac n2\big)}{\Gamma\big(\frac{n-1}{2}\big)} .
\]
Comme $\sqrt K=\frac{\sqrt n\,S}{\sigma}$, on obtient $\E\big(\sqrt n\,S\big)=\dfrac{\sqrt2\,\Gamma\big(\frac n2\big)}{\Gamma\big(\frac{n-1}{2}\big)}\,\sigma$, puis, avec $S_{n-1}=\sqrt{\frac{n}{n-1}}\,S$ :
\[
\E\big(S_{n-1}\big)=\sqrt{\frac{2}{n-1}}\,\frac{\Gamma\big(\frac n2\big)}{\Gamma\big(\frac{n-1}{2}\big)}\,\sigma .
\]
On pose
\[
c_4(n)=\sqrt{\frac{2}{n-1}}\,\frac{\Gamma\big(\frac n2\big)}{\Gamma\big(\frac{n-1}{2}\big)} .
\]
On a alors
\[
\E\left(\frac{S_{n-1}}{c_4(n)}\right)=\sigma ,
\]
et l'on peut donc prendre pour estimateur ponctuel de $\sigma$ : $\hat\sigma=\dfrac{S_{n-1}}{c_4(n)}$.

\begin{rem}
$c_4(n)<1$ : $S_{n-1}$ sous-estime $\sigma$ en moyenne. On le voit aussi sans calcul : $\V(S_{n-1})=\E(S_{n-1}^2)-\E(S_{n-1})^2\ge0$ donne $\E(S_{n-1})\le\sigma$. Pour $n=2$, $\Gamma(1)=1$ et $\Gamma(\frac12)=\sqrt\pi$ donnent $c_4(2)=\sqrt{2/\pi}\approx0{,}798$. Le calcul de $c_4$ utilise la loi du $\chi^2$, donc la normalité.
\end{rem}

\newpage
{\raggedright\Huge\bfseries Estimateur ponctuel de l'écart-type~$\sigma$\\[4pt] d'une population à partir\\[4pt] de l'étendue~$R$\par}

\bigskip
Soient $(X_i)_{1\le i\le n}$ $n$ variables aléatoires indépendantes suivant toutes la loi $\N(\mu;\sigma)$. Les variables centrées réduites $X_i^*=\frac{X_i-\mu}{\sigma}$ ont pour densité $\varphi$ et pour fonction de répartition $\Phi$.

L'étendue $R$ (\emph{range} en anglais) est définie par
\[
R=\max_i(X_i)-\min_i(X_i).
\]
On a
\[
\frac{R}{\sigma}=\frac{\max_i(X_i)-\mu}{\sigma}-\frac{\min_i(X_i)-\mu}{\sigma}.
\]
Considérons
\[
M_n=\frac{\max_i(X_i)-\mu}{\sigma}=\max_i(X_i^*)\qquad\text{et}\qquad m_n=\frac{\min_i(X_i)-\mu}{\sigma}=\min_i(X_i^*).
\]
La fonction de répartition de $M_n$ est donnée, pour tout réel $t$, par
\[
\P(M_n<t)=\P\Big(\bigcap_i\{X_i^*<t\}\Big)=\prod_i\P(X_i^*<t)=\big(\Phi(t)\big)^n ,
\]
où l'indépendance justifie le produit. Ainsi, la densité de $M_n$, notée $f_{M_n}$, est obtenue par dérivation :
\[
f_{M_n}(t)=n\,\varphi(t)\,\big(\Phi(t)\big)^{n-1}.
\]
En outre, $m_n=\min_i(X_i^*)=-\max_i(-X_i^*)$. La loi normale centrée réduite est symétrique : $-X_i^*$ a la même loi que $X_i^*$, donc $\max_i(-X_i^*)$ a la même loi que $M_n$ et $\E(m_n)=-\E(M_n)$. Ainsi
\[
\E\left(\frac R\sigma\right)=\E(M_n)-\E(m_n)=2\,\E(M_n).
\]
On notera $d_2=2\,\E(M_n)$, et l'on a
\[
d_2(n)=2\int_{-\infty}^{+\infty}n\,t\,\varphi(t)\,\big(\Phi(t)\big)^{n-1}\dd t,
\qquad\text{soit}\qquad \E\left(\frac R\sigma\right)=d_2 .
\]
Si l'on dispose de $k$ échantillons \emph{de même taille} $n$, d'étendues $R_1,\dots,R_k$, et $\overline R=\frac1k\sum_jR_j$, alors $\E(\overline R)=d_2\,\sigma$ par linéarité. D'où l'estimateur ponctuel
\[
\hat\sigma=\frac{\overline R}{d_2}.
\]
Il ne mesure que la dispersion \emph{à l'intérieur} des échantillons : décaler toutes les valeurs d'un échantillon ne change pas son étendue. C'est le « $\sigma$ court terme » des calculs de capabilité.

\begin{note}[Pourquoi l'hypothèse de normalité est essentielle]
La symétrie et la forme de $\Phi$ interviennent dans $d_2$. Pour une autre loi, la constante change. Pour $n=5$ : loi normale, $d_2=2{,}326$ ; loi uniforme, $d_2=\sqrt{12}\cdot\frac46\approx2{,}309$ ; loi exponentielle, $d_2=1+\frac12+\frac13+\frac14\approx2{,}083$. Appliquer la table normale à des données exponentielles sous-estime donc $\sigma$ d'environ 10 \%.
\end{note}

\begin{rem}
Pour $n=2$ et $n=3$, on a les valeurs exactes $d_2(2)=\frac{2}{\sqrt\pi}\approx1{,}128$ et $d_2(3)=\frac{3}{\sqrt\pi}\approx1{,}693$. Il n'existe pas de formule générale à l'aide de la fonction $\Gamma$ : $d_2$ se calcule numériquement.
\end{rem}

$d_2$ et $c_4$ se calculent numériquement (voir ci-dessous), ou l'on utilise des tables.

\section*{Exemples de code en R}

Fonction $d_2$, par l'intégrale ci-dessus :
\begin{code}
d2 <- function(n) 2*integrate(function(x) n*x*dnorm(x)*pnorm(x)\^{}(n-1),\\
\phantom{d2 <- function(n) 2*integrate(}-Inf, Inf)\$value
\end{code}
ou en utilisant la loi de l'étendue de $n$ normales, disponible dans R par \texttt{ptukey} avec \mbox{\texttt{df = Inf}} (l'espérance d'une variable positive est l'intégrale de $\P(R>w)$) :
\begin{code}
d2 <- function(n) integrate(function(w)\\
\phantom{d2 <- }1 - ptukey(w, nmeans = n, df = Inf), 0, Inf)\$value
\end{code}
Fonction $c_4$ :
\begin{code}
c4 <- function(n) sqrt(2/(n-1))*gamma(n/2)/gamma((n-1)/2)
\end{code}
On retrouve bien :
\begin{code}
> d2(3)\\
{[1]} 1.692569\\
> c4(3)\\
{[1]} 0.8862269
\end{code}

\begin{note}[Correction par rapport à la version de 2019]
La version de 2019 proposait une seconde fonction \texttt{d2} « en utilisant la fonction Gamma ». Elle calculait en réalité $c_4$ (même formule que la fonction \texttt{C4}). Si on la définissait après la première, \texttt{d2(3)} renvoyait $0{,}886$ au lieu de $1{,}693$. Elle est remplacée ci-dessus par la version \texttt{ptukey}. Les trois fonctions ont été exécutées le 30 octobre 2026.
\end{note}

\end{document}
